\Th
	Обозначим через $d_n$ число многочленов степени $n$, неприводимых над $\Z_p[x]$. Тогда для любого $n \in \N$ выполено следующее равенство:
	$d_n = \frac{1}{n}\sum\limits_{k \mid n}p^k\mu\left(\frac {n}{k}\right)$

\Proof
	Поле $\mathbb{F}_{p^n}$ является полем разложения многочлена $x^{p^n} - x \in \Z_p[x]$. Рассмотрим разложение $x^{p^n} - x = f_1\dotsb f_s$ на неприводимые сомножители и докажем, что $f_1, \dotsc, f_s \in \Z_p[x]$ -- это в точности все различные неприводимые над $\Z_p[x]$ многочлены степеней $k$ таких, что $k \mid n$.
	
	Многочлены $f_1, \dotsc, f_s$ различны потому, что у $x^{p^n} - x$ нет кратных корней. Рассмотрим теперь многочлен $f_i$ степени $k_i$ для произвольного $i \in \{1, \dotsc, s\}$ и его произвольный корень $\alpha \in \mathbb{F}_{p^n}$. Тогда расширение $\Z_p(\alpha)$ имеет степень $k_i$. Но $\Z_p(\alpha) \subset \mathbb{F}_{p^n}$, поэтому $k_i \mid n$. С другой стороны, если $f \in \Z_p[x]$ -- неприводимый многочлен степени $k \mid n$ и $\alpha \in \overline{\Z_p}$ -- его корень, то $|\Z_p(\alpha)| = p^k$, поэтому $\alpha \in \mathbb{F}_{p^k} \subset \mathbb{F}_{p^n}$. Значит, $f = f_i$ для некоторого $i \in \{1, \dotsc, s\}$.
	
	Подсчитывая число корней в левой и правой части равенства $x^{p^n} - x = f_1\dotsb f_s$, получаем, что $p^n = \sum\limits_{k \mid n}k d_k$. Тогда, по формуле обращения Мебиуса, $n d_n = \sum_{k \mid n}p^k\mu(\frac{n}{k})$, что и требовалось.
\EndProof
